\documentclass[12pt]{article}
\usepackage[margin=0.7in]{geometry}

% \usepackage[margin=0.8in]{geometry}

% \usepackage{newpxtext}
% \usepackage{newpxmath}

\usepackage[tt=false]{libertine}

% \usepackage{newtxtext}
% \usepackage{newtxmath}

\usepackage[stretch=10]{microtype}

\linespread{1.05}

\usepackage{enumitem}
\setlist{nosep}

\usepackage[hidelinks, linktoc=all]{hyperref}

% Answers to the review questions. In the PDF this typesets the answer as an
% indented block right after the question; on the web the same environment is
% turned into a click-to-reveal dropdown (see web-pipeline/qa-details.lua).
% Pandoc does not implement \NewDocumentEnvironment, so it leaves the environment
% as a Div the filter can restructure, while LaTeX honours the definition here.
\NewDocumentEnvironment{answer}{}
  {\par\smallskip\noindent\ignorespaces}
  {\par\smallskip}

\title{Ethics}
\date{}
\author{}

\begin{document}

\maketitle
\tableofcontents
\newpage

\section{Metaethics}
The study of the foundation of morality.

\noindent Theme: Debate over whether moral facts exist universally and independently of human minds.
Is there a single standard of Right and Wrong that applies to human beings everywhere?

\subsection{Cognitivism vs Non-cognitivism}
Are moral statements like "stealing is wrong" stating a fact about the world,
or just expressing one's personal feelings?

\textbf{Cognitivism} This view believes that moral statements are truth apt,
i.e. they can be objectively True or False. Cognitivists have to answer "\textit{what exactly makes these statements true?}".

\textbf{Non-cognitivism} This view argues that moral statements are \textbf{not} capable of 
being True or False. Instead, they are expressions of our inner emotions, attitudes, or personal preferences.
Note that being a non-cognitivist does not make one a moral skeptic (rejecting morality altogether).

\subsection{Emotivism}
Emotivism is a specific type of non-cognitivism, which claims that moral statements are expressions of emotions or feelings.
"\textit{Stealing is wrong}" is expressing our emotion of disapproval.

\subsection{Error Theory \& Extreme Nihilism}

\textbf{Error theory} is a form of cognitivism. It agrees moral statements \textit{attempt} to express objective facts.
But error theorists claim that all ethical convictions are false beliefs.
For the statement "\textit{Stealing is wrong}", Error Theory argues this is making a claim about a property ("wrongness")
that does not exist in the world. Hence, any moral statement is technically a false belief.

\textbf{Extreme Nihilism} extends Error theory to everyday lives. It claims that nothing actually has value, i.e. the concepts of right vs wrong,
good vs bad, simply do not exist. The entire concept of morality is an illusion.
The argument:
\begin{enumerate}
    \item We should only believe in something (or accept we have evidence for it) if that thing helps us explain the physical world we observe around us.
    For example, we believe in gravity because it explains why an apple falls to the ground.
    \item "Moral hypotheses" (like the idea that an action was "evil" or "unjust") never actually help us explain why anything happens in the physical world. 
    Because moral facts don't explain physical observations, we have absolutely no evidence that our moral opinions are true.
\end{enumerate}
This means that there are no moral constraints in the world (right and wrong are illusions), so we can do anything we want (no rules governing our behaviour).

\subsection{Moral Realism}
Moral realism is a form of cognitivism.

\noindent \textbf{Claim}: There are objective, mind-independent moral facts.
Mind-independent means that the moral facts remain objectively and completely true, similar to math or physical facts,
regardless of what some individuals or societies think or feel.
The statement "\textit{Stealing is wrong}" is a factual, universal truth regardless of what some individuals may think.

\noindent Why might one be a moral realist:
\begin{enumerate}
    \item \textbf{Establish secular objectivity} Provides a basis for moral truths that exist regardless of cultural or personal beliefs.
    For example, "\textit{Murder is wrong}" is a natural fact of the universe that we can discover, just like gravity or biology.
    \item \textbf{Validate Human Intuition} Moral realism allows us to do justice to our innate sense of right and wrong.
    It validates our strong intuition that ethical issues are genuine facts, rather than mere matters of opinion.
    For example, "\textit{Torturing an innocent puppy for fun is wrong}", moral realism would validate our feelings that that act is factually, objectively wrong.
\end{enumerate}

\subsection{Ethical Naturalism}
A version of moral realism (right and wrong are objective facts).
Claims:
\begin{enumerate}
    \item \textbf{Moral realism} There are objective, mind-independent moral facts.
    \item \textbf{Metaphysical Naturalism} The belief that all facts are natural facts.
    (Since all substances and properties that exist are purely natural, then nothing supernatural exists).
    Moral properties are strictly defined by facts about the natural, physical world.
    The "wrongness" of murder is just a complex but natural fact, perhaps rooted in biological survival etc.
    \item \textbf{Epistemic Naturalism} Claims that because moral facts are just natural facts,
    we can discover and verify them using the same methods in the natural sciences like physics.
    For example, to know if something is "morally good", we may observe its physical and psychological effects on humans.
\end{enumerate}
Objective right and wrong exists, and are completely natural, physical facts, therefore we can use science, observation, logic to figure out what those moral facts are.

\vspace{1em}
\noindent\textbf{Example of Ethical Naturalism} Hedonistic Utilitarianism: It claims that 'Goodness' is exactly equal to 'Pleasure.'
Because pleasure is a biological, psychological, and entirely natural phenomenon that can be scientifically observed, this theory successfully grounds a moral concept in the physical, natural world.

\subsection{Ethical Supernaturalism}
Is a form of moral realism, agreeing that there are objective, mind-independent moral facts.
However, ethical supernaturalists claim that ethical propositions are grounded in a divine or supernatural realm, rather than in nature.

\textbf{Divine Command Theory} Some supreme being, perhaps God, is the some ultimate source of morality. An action is morally good strictly because it is commanded by a divine being.


\subsection{Two Objections Against Ethical Naturalism}
\subsubsection{Science vs Ethics (Harman's Challenge)}
Recall that ethical naturalism claims that we can use scientific methods to discover moral facts.
Gilbert Harman's counter-argument asserted that moral facts are unnecessary to explain our observations.

\noindent A scientific example: When a physicist looks into a cloud chamber and observes a proton vapour trail,
the existence of the physical proton is strictly necessary to explain that observation.

\noindent A moral example: When we see someone stealing a homeless person's wallet,
we observe that "\textit{that's wrong}". The only reason for this observation is our own psychology, upbringing etc.
Note that we did not anything physical, objective property called "wrongness" to explain that observation.

\noindent Science relies on Ockham's Razor: we should only believe in things that are strictly necessary to explain our observations.
Objective, moral facts are not necessary to explain our moral observations.
Since ethical naturalism uses scientific standards to claim the existence of moral facts,
the same scientific standards disproved the existence of moral facts.

\textbf{Counter to Harman} Mathematical proofs and equations are widely accepted
as objectively true, although they do not require empirical, physical observation. 
E.g. in equations like 5 + 7 = 12, We can't physically perceive or touch the number 12.
In the same way, it is impossible to physically perceive abstract moral qualities.
Therefore, we should accept moral facts just as we accept something like mathematics,
as not all objective facts have to be physical things we can observe. 

\subsubsection{Moore’s Open Question Argument}
Ethical naturalism claims that moral concepts like "Goodness" are just natural properties like "pleasure" or "biological survival".
G.E. Moore argues that identifying a moral property with a natural property is a fallacy.

\noindent A closed question is a tautology, like "\textit{Is a bachelor an unmarried man?}".

\noindent An open question implies that the answer is not contained within the definitions of the terms alone.
For any proposed natural definition of "good", (like "pleasure"), the question
"\textit{Is pleasure good?}" would be a tautology under ethical naturalism.
However, asking if something is good is always an open, meaningful question.
Therefore, Moore argues that "Goodness" is an unanalysable, non-natural property that cannot be defined by anything else. 

\noindent\textbf{Counter to Moore} Meaning is not the same as synthetic identity.
Moore's argument relies on analytic identity, that for two things to be identical,
their words must have the exact same dictionary meaning.
However, there is also synthetic identity, like "\textit{Water is H20}", 
where it might have been an open question hundreds of years ago, but after scientic discovery, the two concepts are actually the same thing in the real world,
even if their dictionary definitions are different.

\noindent Hence, questions like "\textit{Is pleasure good?}" may be an open question today,
but it may mean that we are still in the process of scientifically discovering the synthetic identity of "Goodness". Moral properties could be natural properties even if their concepts are different.

\subsection{Reductive Naturalism}
Reductive naturalism presents a solution to Harman's challenge that we cannot physically observe moral facts.
Reductive naturalism argues that we can construct moral facts out of other natural, observable facts.
Scientific example: Colour can be broken down into a combination of facts about light waves, physical surfaces and human eyesight.
Similarly, "goodness" is simply a combination of biological, psychological and social facts.
Reductive naturalists use Functionalism to do this, which is the idea that something is good if it effectively fulfills its purpose.
For example, a watch is "good" if we can observe that it keeps accurate time, if it's durable, if it fits well, etc.
We can apply this to human beings, where a "good" human action might be one that adequately fulfills the function of promoting human survival,
societal harmony, or psychological well-being. 

\noindent\textbf{The problem with Functionalism} is that humans do not have a single, objective, universally agreed-upon "purpose" or "function".
For example, if a naturalist claims a human's purely natural function is "biological reproduction and survival," does that mean a person who chooses to be childless or sacrifices their life for a stranger is functionally "broken" or "morally bad"?
Reducing human morality to biological functions often leads to highly controversial and counterintuitive conclusions.

\newpage

\section{Moral Relativism}
The meta-ethical theory that no absolute, universal or objective moral principles exist.

\noindent Theme: Why do different groups of humans have completely different moral facts?

\vspace{1em}
\noindent Note: Moral Relativism is the overarching theory that no universal truths exist. Cultural Relativism is a specific type of moral relativism where the 'framework' defining truth is a society/culture.
Subjectivism is another type where the 'framework' is the individual.

\subsection{Descriptive vs Normative Claims}
Descriptive claim - a statement of empirical fact, describing what is currently happening without passing judgement.

\noindent Normative claim - A statement of value, makes a judgement about how the world ought to be.

\noindent Is-ought fallacy (Hume's Law): One cannot jump from a descriptive statement to a normative conclusion.
Just because something is a certain way, doesn't mean that it ought to be that way.

\subsection{Cultural Relativism (Descriptive vs Normative)}
Descriptive cultural relativism - purely observational, empirical claim about human societies. What is morally right and wrong
varies between cultures. Crucially, what varies is what people in different cultures think or believe is right and wrong.

\noindent Example: "\textit{In some cultures, people think gay marriage is okay; in others they don't}".

\noindent Example: "\textit{In the 1800s, the culture of American South believed slavery was morally acceptable. Today, modern American culture believes slavery is deeply immoral}".

Normative cultural relativism - What is morally right and wrong varies between cultures, 
and what varies is what \textit{really is} morally right and wrong.

\noindent Example: "\textit{In some cultures gay marriage really is permissible; in others it isn't}".

\noindent Example: "\textit{In the 1800s, slavery was actually morally right for the American South because their culture approved of it. Today, it is actually morally wrong for us}".

\noindent Consequences of Normative Cultural relativism
\begin{enumerate}
    \item We cannot criticize or judge other cultures.
    If a culture's own code is the ultimate standard of right and wrong, we have no right to interfere with them.
    \item We cannot criticize our own culture.
    If the majority of your culture decides what is right, then anyone who disagrees with the majority is, by definition, acting immorally.
    Under this theory, social reformers like Martin Luther King Jr. or Mahatma Gandhi were technically acting immorally at the time, because they were protesting against the established moral code of their own cultures.
    \item "Moral Progress" is a myth. We may think that humanity has improved over time, e.g. abolishing slavery.
    But according to Normative Cultural Relativism, there is no objective standard to move towards, and any change is just a shift in cultural tastes and norms.
\end{enumerate}

\subsection{Descriptive vs Metaethical vs Normative Moral Relativism}
Descriptive Moral Relativism - a broader claim about humanity, that empirically, deep and widespread moral disagreements exist across different societies, 
and crucially, these disagreements are much more significant than whatever moral agreements there might be.

\noindent Metaethical Moral Relativism - Asserts that the truth or falsity of a moral judgement is not absolute or universal. (No objective truth exists.)
The truth of a moral claim is entirely relative to the framework (culture/group).

\noindent Normative Moral Relativism - extension of the metaethical view, adding in normative implications on what we should do.
Because there are no universal moral truths, the correct moral action is always dependent on context (often culture, but can also be individual).
Further prescribes that we ought not to judge the moral beliefs/practices of other cultures, as no single moral code is superior.

\noindent\textbf{Example}
\begin{itemize}
    \item Descriptive: Anthropologists observe that Culture A eats cows, while Culture B worships them.
    \item Metaethical: The statement 'Eating cows is right' is true for Culture A and false for Culture B, and neither culture is objectively 'more correct' than the other.
    \item Normative: Because truth is relative, Culture A ought not to judge or force its dietary laws onto Culture B.
\end{itemize}

\subsection{The Cultural Difference Argument}
\begin{enumerate}
    \item Premise: Different societies have different moral codes.
    \item Conclusion: There is no objective truth in morality, and we cannot judge the moral codes of other societies.
\end{enumerate}
Logical fallacy: This arugment violates Hume's Law, i.e. it jumps from a descriptive claim to normative claim.
Just because cultures disagree about a moral issue, it is logically invalid to conclude that a universal truth ought not or does not exist.
Similar to disagreements about scientific facts, one culture may be correct while another is mistaken.
The fact that cultures disagree on moral issues only proves that figuring out moral truth is difficult, but does not logically prove that an objective moral truth does not exist.

\noindent\textbf{Example}
\begin{enumerate}
    \item In the 18th century, some cultures believed that slavery was morally acceptable, while other cultures believed it was a horrific moral evil.
    \item Therefore, there is no objective moral truth about whether slavery is right or wrong. 
    It is purely a matter of cultural opinion, and we cannot say that one culture was objectively "better" or "more correct" than the other.
\end{enumerate}
The mere existence of a disagrement doesn't prove that a right answer doesn't exist. It is entirely possible that 
an objective moral truth does exist, and the cultures that endorsed slavery were simply horribly mistaken.

\subsection{Problem with Moral Relativism and Subjectivism}
% e.g., collapse into non-cognitivism?
Problems with applying moral relativism:
\begin{enumerate}
    \item \textbf{Boundaries} If a culture decides what is morally true, how to define the boundaries of that culture?
    Is it a country? Or a state?
    \item \textbf{Reformer's Dilemma} If a culture's current standards are the definition of "right",
    then anyone trying to change those rules (like reformers fighting to end slavery) are technically acting immorally.
\end{enumerate}
\textbf{Slippery slope to subjectivism} Boundaries of a culture are hard to define, so relativism might collapse into subjectivism.
If a dominant culture can have its own truth, why can't a sub-culture? If a sub-culture can, why can't a specific family? And if a family can, why can't a single individual?

\noindent Subjectivism argues that right and wrong are based on individual feelings, attitudes, or opinions,
i.e. morality is entirely up to the individual person.
Critics argue that moral relativism tends to slide all the way down to subjectivism, and then ethics doesn't give any meaningful set of rules.

\noindent\textbf{Collapse into non-cognitivism}
Relativists say truth depends on the culture. Subjectivists say truth depends on the individual.
They share the same problem: if moral 'truth' is just whatever a person/group currently approves of, then a moral statement loses its factual weight.
Moral statements are just you venting your personal feelings, and have no actual truth value.
The theory collapses completely into Non-Cognitivism.
(Recall that non-cognitivism means moral statements are not capable of being True or False,
instead they are just expressions of emotions.)

\noindent Saying 'Stealing is wrong' is no longer a cognitive fact about the world; it simply translates to 'I personally dislike stealing!' Because it just becomes a way to vent personal feelings rather than state facts.

\noindent\textbf{Harman's defense of Moral Relativism} 
Harman argues that there are indeed no universal moral truths, 
but that ethics can still be treated as a system of rigorous, logical facts.
\begin{enumerate}
    \item \textbf{Frames of reference} When throwing a ball on a moving train, the velocity of the ball is relative to a specific frame of reference.
    Harman argued that morality works the same way, i.e. "Stealing is wrong" should actually mean "Stealing is wrong \textit{relative} to moral framework X".
    \item \textbf{Implicit agreements} Moral truths are relative to a framework, and that framework is an implicit agreement among a group of people.
    This agreement determines what actions are deemed acceptable or "wrong" within that specific context.
    \item \textbf{Inner Judgements} Because morality is based on these shared agreements,
    it only makes logical sense to say someone "ought" or "ought not" to do something if we assume they are actually plugged into our moral framework and share our basic motivations.
\end{enumerate}
By defining moral statements as relative to an implicit agreement, this saves ethics from collapsing into meaningless emotion.
If someone says, "Stealing is wrong", they are actually making a factual, cognitive claim about the relationship between an action and a specific framework.
The truth value of the statement then depends on the framework (the implicit soceital agreement).

\newpage

\section{Utilitarianism}
\subsection{Consequentialism}
The moral status of an action depends entirely on their producing the best consequence.
The implication is that actions themselves are not inherently "good" or "evil"; their moral worth is judged purely by the results they produce in the real world.

\vspace{1em}
\noindent\textbf{Difference between consequentialism and utilitarianism}
Utilitarianism is just one specific type of consequentialism.
Consequentialism just says 'maximize the good consequences'.
It doesn't define what 'good' is.
A consequentialist could believe the ultimate 'good' is Knowledge, Beauty, or Equality.
Utilitarianism specifically defines the 'good' as Well-being/Utility.
Therefore, all utilitarians are consequentialists, but not all consequentialists are utilitarians.

\subsection{Utility Principle}
Always act to maximise "the greatest good for the greatest number".
An action is morally right if it produces the greatest net balance of well-being over suffering for everyone affected.

\noindent \textbf{Impartiality} - Everyone's well-being counts exactly the same.

\vspace{1em}
\textbf{Note}: The Utility Principle does not necessarily mean we must always choose the action that benefits the greatest number of people.
It requires us to maximize the total net utility.
For example, an action that gives a massive amount of happiness to 2 people might be morally superior to an action that gives a tiny, barely noticeable amount of happiness to 100 people.
We calculate total overall well-being, not just a head-count.

\subsection{Hedonistic Utilitarianism}
A specific form of utilitarianism. Hedonism (Bentham and Mill) says that pleasure is the only intrinsic good and pain is the only intrinsic evil.
Hence, hedonistic utilitarianism says that the morally right action is one that maximises the total net amount of pleasure and minimises the total net amount of pain.

\noindent Similarly, it is impartial - taking equal consideration of pleasure/pain of all affected beings.

\noindent\textbf{Recall} that Hedonistic Utilitarianism is an example of Ethical Naturalism.

\subsection{Actual vs Foreseeable Consequence}
Actual consequence utilitarianism evaluates an action's morality based on its actual, final outcomes.
Foreseeable consequence (or expected) utilitarianism judges actions based on predicted, reasonable expectations at the time of the decision.

\noindent Scenario: Imagine a person in 1938 saves someone from drowning, however the person they saved turned out to be Adolf Hitler.
Actual consequence utilitarianism says saving the drowning man was a morally wrong action. The actual consequence of the action was a massive net negative for global well-being.
However, this makes moral decision-making nearly impossible. Because you cannot see the future, you can never be 100\% certain you are doing the right thing at the time you act. 
You can only evaluate if an action was right or wrong retrospectively, sometimes years later.

\noindent Forseeable consequence utilitarianism says saving the drowning man was the morally right action.
Based on the information available to a reasonable person at the time, saving a life has a massively high expected utility. The fact that he later became a dictator was an unforeseeable statistical anomaly.
This version uses foresight instead of hindsight. It allows human beings to make moral decisions using probability and reasonable expectations, meaning you can be confident you are doing the right thing at the moment you act, even if terrible luck ruins the actual outcome later.

\subsection{Benthamic Principle: ``Can they suffer?'' \& Moral Community}
Moral community is a group of beings that are recognized as having moral status and worth, and whose interests we are morally required to consider when we make decisions.
Before Bentham, the boundaries of the moral community was drawn around Rationality, Language or Species Membership (being human).
Bentham said that under hedonistic utilitarianism (only pain and pleasure matter), reason and language are completely irrelavant.

\noindent He famously wrote: "The question is not, Can they reason? nor, Can they talk? but, Can they suffer?"
If the capacity to suffer (sentience) is the only requirement for moral standing, the boundary of the moral community explodes outward.
Under this principle, a unit of pain felt by a dog, a pig, or a cow counts exactly the same in the utilitarian cost-benefit calculation as a unit of pain felt by a human being.
To a utilitarian, if an action causes massive suffering to animals for only a minor amount of human pleasure (like eating a specific type of meat), that action is mathematically and morally wrong.

\noindent\textbf{Problem with Bentham}
If suffering is the only metric, how do we accurately compare human suffering to animal suffering?
Does a pig's physical pain count exactly the same as a human's complex psychological grief?
It makes the utilitarian calculation incredibly difficult.

\subsection{Argument from Marginal Cases}
Marginal human beings refer to humans whose cognitive abilities are at the margins of typical human capacity, like infants, or individuals with severe cognitive disabilities.
\begin{enumerate}
    \item \textbf{Premise}: If it is immoral to kill and eat "marginal" human beings, and to painfully experiment on them, then it is also immoral to treat non-human animals this way. (Because an adult chimpanzee or a pig often has the exact same, or even higher, cognitive capacity and ability to feel pain as a human infant).
    \item \textbf{Premise}: It is (almost) always immoral to kill and eat "marginal" human beings and to painfully experiment on them. (This is a basic moral intuition almost everyone shares).
    \item \textbf{Conclusion}: It is (almost) always immoral to kill and eat animals, and to painfully experiment on them.
\end{enumerate}
Speciesism Counterargument: The AMC asserts that species membership is not a morally relevant characteristic; treating humans with lower capabilities differently than animals with similar capabilities is a form of prejudice similar to racism or sexism.

\noindent\textbf{Problem with Marginal Cases}
Critics argue that moral status isn't based on an individual's specific cognitive capacity at a given moment, but on what is normal for their species.
Because higher-level rationality is normal for the human species, all humans get moral rights, even if a specific infant or disabled individual currently lacks it.

\subsection{Obligatory vs Supererogatory Actions}
Obligatory actions - Something that you are morally required to do. It would be morally wrong not to do it. (For example, feeding your own children is obligatory). Note: If an action is obligatory, it is automatically permissible, but not vice versa.

\noindent Supererogatory actions - Something that is morally positive and good, but not obligatory. This is an action that goes "above and beyond the call of duty." If you do it, you are a hero, but if you don't do it, you haven't done anything wrong. (For example, running into a burning building to save a stranger, or donating 70\% of your income to charity).

\noindent This distinction collapses under utilitarianism because any action that maximizes overall utility is a strict moral obligation, not optional.
An action is either the single best possible choice (making it obligatory), or it isn't (making it wrong).
Supererogatory actions that maximise utility are thus obligatory.

\subsection{3 Objections Against Utilitarianism}
\subsubsection{Measurement}
If we want to maximise the net balance of well-being over suffering,
how to actually measure, compare and sum up this "utility"?
\begin{enumerate}
    \item \textbf{Premise 1}: If utilitarianism is true, then there must be a precise unit of measurement that determines the amount of well-being that results from an action.
    \item \textbf{Premise 2}: There is no such precise unit of measurement.
    \item \textbf{Conclusion}: Therefore, utilitarianism is false (or at least practically useless for moral problem-solving).
\end{enumerate}
Utilitarians reject Premise 1, arguing that we do not need a perfectly precise, mathematical unit of measurement to make valid moral choices.
\begin{enumerate}
    \item \textbf{Psychological Similarity}: Because human beings share very similar biology and psychology, we can safely estimate others' wants, needs, and what makes them happy.
    \item \textbf{Accumulated Wisdom}: Over thousands of years of human history, we have gathered enough general knowledge to know what outcomes generally promote well-being and what outcomes destroy it.
\end{enumerate}

\subsubsection{Deliberation}
Utilitarianism is impractical as a decision-making guide because calculating the total utility of every possible action is impossible, time-consuming, and often counterintuitive.
Humans are not "moral calculators" capable of crunching infinite data points in real-time.
We do not have perfect knowledge of the future, so we simply don't know how things will actually turn out.

\noindent Utilitarians split moral action into two distinct concepts
\begin{enumerate}
    \item \textbf{Standard of Rightness} The utilitarian formula itself that defines what makes an action morally correct.
    \item \textbf{Decision Procedure} The practical method used to make choices. Relying on rules of thumb (like "don't lie", "keep your promises", or "don't steal") is the best decision procedure.
\end{enumerate}
\textbf{Basketball analogy} Basketball player does not need to pause to calculate the precise physics equations of velocity, gravity, and wind resistance.
They just use their practiced habits and intuition to catch the ball. Similarly, you don't need to calculate the equations of morality every time you act; 
you just follow the rules of thumb that have been proven to maximize well-being.

\subsubsection{Undermining Trust}
Under utilitarianism, no action is intrinsically wrong. Stealing, lying, breaking promises, killing etc are evaluated strictly by their consequences.
Since trust requires that people consistently follow rules (e.g., honesty, fairness), then utilitarianism's focus on the "best outcome" in each situation makes behavior unpredictable, undermining trust.

\noindent\textbf{Transplant Example} If doctors might kill one healthy patient to save five with their organs, people would fear visiting hospitals, destroying trust in the medical system.

\noindent\textbf{Justice Example} If a judge frames an innocent person to stop a riot, the long-term trust in the law as a fair institution is destroyed.

Utilitarians might respond by arguing that social trust has massive utility.
In the transplant example, if public trust is eroded, people stop going to hospitals, and thousands of people die of preventable diseases.
The collapse of trust creates a massive amount of long-term suffering, so the utilitarian calculation actually forbids the doctor from killing you.

\subsection{Agent-Relative vs Agent-Neutral Reasons}
Agent-neutral reasons are universal, holding the same value for everyone regardless of who acts (e.g., "maximize overall happiness").
Utilitarianim is agent-neutral.

\noindent Agent-relative reasons depend on the individual's specific circumstances or relationships.
Agent-relative reasons are not just subjective, personal feelings (e.g., 'I only help people I think are attractive').
They are still universal moral rules, but they are applied to specific relationships.
For example, 'Every parent has a duty to save their own child first.'
This is a universal rule (it applies to all parents everywhere), but it is agent-relative because the specific child you are obligated to save depends on who you are.

\noindent\textbf{Transplant Example} Under this framework, the doctor can say, "I have a specific, agent-relative constraint against me personally committing a murder, even if my murder prevents five other deaths."

\newpage 

\section{Deontology}
Deontology is a duty-based ethical theory asserting that the morality of an action is based on whether the action itself is right or wrong under a series of rules, rather than based on the consequences.

\subsection{Categorical Imperative}
According to Kant, the Categorical Imperative (CI) is the central, unconditional moral law that dictates actions must be based on duty rather than desire.

\noindent\textbf{Hypothetical Imperative (Desire-based)}
"\textit{If you want X, do Y}". These rules are based entirely on our personal inclinations or desires.

\noindent In contrast, CIs are \textbf{reason-based}.
"\textit{I should do X whether or not I want anything else}".
These are moral commands based on reason alone, completely ignoring what you personally want.
In other words, the rule is followed purely out of a sense of duty, rather than a desire for a reward, this action holds genuine moral goodness.

\subsection{Principle of Universalisability}
This is the first formulation of Kant's Categorical Imperative.

\noindent PoU says "you should act only on principles that you would be willing to have absolutely everyone else act on as well".
In other words, an action is morally right if and only if it is universalisable.

\noindent\textbf{The Universalisability Test}
\begin{enumerate}
    \item Imagine a hypothetical world in which absolutely everyone performs the action you are about to take.
    \item Confirm that in this hypothetical world, both of the following statements are true:
    \begin{itemize}
        \item The goal of the action can still actually be achieved in such a world.
        \item Nothing essential to the agent’s (your) will would be endangered if everyone in the world performed this action.
    \end{itemize}
\end{enumerate}
If both of those statements are true, the action is universalizable, which means it is morally right. 
If either statement is false, the action is a contradiction and is morally wrong.

\noindent The fundamental thought behind this principle is that acting morally strictly requires being impartial, and not making an exception for themselves.

\noindent\textbf{Example - Lying}

\noindent Imagine there is something very expensive you want to buy, but you lack the money.
You have a wealthy acquaintance who you know will give you a loan if you simply ask.
This loan would fully cover the cost of the item you want.
You know for a fact that if you never repaid the loan, you would not suffer any punishment or retaliation from this acquaintance or anyone else.

\noindent Applying the test, this action is immoral (fails statement 1).
If absolutely everyone in the world constantly made false promises to get loans, the entire social concept of a "promise" or a "loan" would completely collapse. No one would ever believe anyone else, and consequently, no one would ever lend out money.
Because your goal (getting the loan by making a promise) is logically impossible to achieve in a world where everyone acts the same way, the action contradicts itself.

\noindent\textbf{Example - Not helping}

\noindent Imagine this maxim: "I will never help anyone, and I expect no one to ever help me."
Applying the test, this world can technically exist (passes statement 1), although it would be a cold, cruel world.

\noindent However, Kant points out a biological and rational fact about human beings: we are not self-sufficient.
As a rational human being, you inevitably will that somebody helps you when you are in danger.
This contradicts the will that nobody helps you (fails statement 2).
Therefore, we have a duty to help others.

\vspace{1em}
\noindent\textbf{Problem with Universalisability Test}
The Universalizability test is highly vulnerable to how you phrase your maxim, leading to "False Positives.
If your maxim is "I will rob a bank to get money," it fails the test (if everyone robs banks, the concept of money collapses).
But if you phrase it very specifically: "I will rob a bank only on Tuesdays in July while wearing a silly hat", it might actually pass the test.
If everyone in the world who wore a silly hat on a Tuesday in July robbed a bank, the banking system would likely survive.
The test can be tricked by overly specific wording.

\subsection{Principle of Humanity}
This is the second formulation of Kant's Categorical Imperative.
An action is morally right if and only if it treats human beings as ends in themselves, and never merely as means to an end.
\begin{enumerate}
    \item Never use others as \textbf{mere} means to an end.
    To use someone as a "mere means" is to use them strictly as a tool for purposes that are not their own purposes.
    For example, if you make a lying promise to get a loan, you are using that person purely as a financial tool.
    However, Kant notes that it is morally permissible to use others for purposes that they actively share with you  (e.g., you use a barista to get coffee, but they share the purpose of completing that transaction).
    \item Always treat others as ends in themselves.
    To treat someone as an "end in themselves" means treating them with respect and recognizing them as a rational, autonomous being with intrinsic value and dignity.
    For example, telling the truth allows others to use their own rationality to decide if they want to help you.
\end{enumerate}
"Human beings" refer to all rational, autonomous agents.
\begin{itemize}
    \item Rationality: The power to reason and to use that reason to govern your own life.
    \item Autonomy: The capacity to govern your own life by making your own free choices.
\end{itemize}
All rational, autonomous beings are intrinsically valuable, and they all possess equal intrinsic value.

\noindent\textbf{Implications on our moral obligations}
\begin{itemize}
    \item Negative Obligations (Constraints):
    We are morally required to refrain from diminishing other people's rationality and autonomy, as well as our own. 
    This creates strict obligations not to lie, steal, enslave, rape, or murder, because all of these actions involve using people for your own purposes.
    \item Positive Obligations:
    We are also required to actively promote people's rationality and autonomy. 
    For example, you have a moral obligation to perform easy rescues (like saving a drowning toddler) because death destroys a rational autonomous being. 
    It also means we have an obligation to help reduce hunger, great poverty, and powerlessness.
\end{itemize}

\subsection{Perfect vs Imperfect Duties}
Kant divided our moral obligations (generated by the CI)
into Perfect and Imperfect duties.

\noindent\textbf{Perfect duties} are mandatory, strict rules that must always be followed with no exceptions.
Violating them creates a "contradiction in conception" (Statement 1, logically impossible to universalize).
They are usually negative duties:
\begin{itemize}
    \item Do not lie
    \item Do not steal
    \item Do not murder
    \item Do not break promises
\end{itemize}

\noindent\textbf{Imperfect duties} are meritorious, general moral obligations we should adopt as lifelong goals,
but you have the flexibility to choose when, how, and toward whom you fulfill them.
Violating them is not a logical contradiction, but a "contradiction in will" (Statement 2, we would not want to live in a world where everyone acted that way).
They are usually positive duties:
\begin{itemize}
    \item To help others in need
    \item To develop one's own talents
\end{itemize}

\noindent\textbf{Conflict} If a perfect duty conflicts with an imperfect duty,
we should prioritise the perfect duty because it is a strict obligation.
For example, one must not steal (perfect) even to help someone in need (imperfect).
The ends never justify the means.

Because Perfect Duties are negative constraints that protect basic human autonomy (Contradiction in Conception), while
Imperfect Duties are positive goals (Contradiction in Will).
You cannot actively violate someone's basic autonomy (break a perfect duty) just to achieve a positive goal (fulfill an imperfect duty).

\vspace{1em}
\noindent\textbf{Perfect/Imperfect vs Obligatory/Supererogatory}
Imperfect duties (like helping others or developing talents) are not to be confused with supererogatory actions (like donating large sum to charity).
An imperfect duty is still a strict moral requirement. You must adopt the goal of helping people. The only "imperfect" part is that you have the flexibility to choose when, where, and how you fulfill it.

\subsection{2 Objections Against Deontology}
\subsubsection{Inquiring Murderer}
\textbf{Scenario} Imagine someone you love is fleeing from a murderer and tells you where she is hiding. The murderer comes to your door and inquires about her location. Should you lie?

\noindent Assumptions:
\begin{itemize}
    \item The inquiring murderer is no threat to you personally; he will not attack you if he finds out you lied to him.
    \item If you don’t lie, the murderer will most probably find and kill the intended victim.
    \item If you do lie, the murderer will most probably not find the intended victim.
\end{itemize}

\noindent According to Kant's CI, we should never lie, as lying violates the CI.

\noindent Benjamin Constant's claim is that it would not be immoral to lie to the inquiring murderer. In fact, lying here may be morally required
to save a life.

\noindent\textbf{Kantian defence} Although the CI is a moral absolute,
it may be more flexible than even Kant himself realized.
The CI can demonstrate that it would be morally right to lie in some situations.
Specifically, they argue that lying to the murderer can actually pass the Universalizability test.
\begin{itemize}
    \item The goal of this action can be achieved in this world. 
    In such a world, an inquiring murderer wouldn’t know for sure whether anyone was lying to him. 
    For all he knows, they really have no idea. 
    Contrast this with a world in which everyone makes a lying promise to repay a debt; in that world, as soon as someone promises to repay a debt, you know for sure they are lying.
    \item Nothing essential to the agent’s will would be endangered. 
    In fact, it is protected because more innocent lives would be saved. 
    Your will could enjoy the protection of others who would lie to save you from being killed.
\end{itemize}


\subsubsection{Animal Rights}
Kant's Principle of Humanity specifically states that we must treat beings with rationality and autonomy as ends-in-themselves.
However, this excludes beings lacking rationality, like infants, individuals with severe mental disabilities, and non-human animals from receiving direct moral consideration.

\noindent This is in contrast from utilitarian approach that asks "\textit{Can they suffer?}".

\noindent For example, the statement "\textit{animal cruelty is perfectly fine}" goes against basic human intuition.
Kant created a workaround: He argued that we do not have direct duties to animals, but rather indirect duties to them because of how our actions affect us.
\begin{enumerate}
    \item \textbf{Premise 1}: Cruelty to animals makes you more likely to be cruel toward rational, autonomous beings (other humans). Kant believed that abusing animals degrades our natural sense of empathy and makes us more cruel overall.
    \item \textbf{Premise 2}: It is immoral to act in a way that makes you more likely to be cruel toward rational, autonomous beings.
    \item \textbf{Conclusion}: Therefore, cruelty to animals is immoral.
\end{enumerate}

\noindent Ultimately, our duties toward non-rational animals and nature are really just duties to respect our own humanity, which requires us to cultivate virtues like kindness and gratitude.

\subsection{Agent-Centered vs Victim-Centered Deontology}
\textbf{Agent-centered Deontology} focuses on the moral agent’s actions and intentions.
Agents have special, "agent-relative" obligations to avoid doing wrong, even if doing so results in more harm overall.
For example, parents have agent-relative obligations to their own children. Furthermore, parents are granted the agent-relative permission to save their own children, even if it comes at the expense of other children not being saved.

\noindent Doctrine of double effect: an action with both a good effect and a foreseen, unintended bad side effect may be morally permissible if the intention is solely to achieve the good result.
For example, in the Trolley Problem, if your explicit intention is to save the five people (and the death of the one person on the other track is just a tragic, foreseen side effect), switching the tracks may be permissible. 
However, if your intention is to actively harm the innocent person, it is not permissible.

\noindent\textbf{Victim-centered Deontology} focuses on the rights of the victim, and rules are designed to protect individuals from being used as a mere means to an end.
It is rights-based rather than duty-based.

\noindent For example, in the transplant case, It is wrong to harvest one person's organs to save five, because this directly violates that person's right not to be used, regardless of the overall outcome.
In the Fat Man trolley problem, because pushing the Fat Man off the bridge uses his body as a physical barrier to stop the trolley without his consent, it is strictly forbidden.

\subsection{Paradox of Deontology}
Victim-centered deontology claims that the foundation of morality is the rights of the victim (e.g., the right not to be killed, or the right not to be used as a mere means).
Can this theory be agent-neutral, i.e. everyone shares the exact same universal goal?

\noindent\textbf{Scenario}: A villain is about to murder five innocent people. You can stop him, but only if you murder one innocent person yourself.
The universal goal here would be "Minimize the total number of rights violations in the world".
An agent-neutral view says you must murder the one innocent person (1 violation is better than 5 violations).
The problem is that Deontology strictly forbids you from murdering an innocent person to achieve a good consequence. Therefore, Victim-Centered Deontology cannot be agent-neutral.

\noindent Victim-Centered Deontology must be Agent-Relative.
The personal goal would be "Ensure that I personally do not violate anyone's rights, regardless of what happens around me".
So we do not kill the one innocent person, and keep our hands clean at the expense of five innocent people dying.

\noindent The paradox is that a Victim-Centered theory was supposed to protect the victim's right not to be killed.
We protected the one victim from us, but allowed five other victims to die, just so that we don't break the rules.

\noindent Victim-Centered Deontology accidentally collapses into Agent-Centered Deontology.
It claims to care about the victim, but when push comes to shove, it actually prioritizes the moral purity of the agent over the lives of the victims.
The paradox happens because Victim-Centered theories enforce an Agent-Relative constraint: 'I must not violate a right.'
By focusing on keeping my hands clean, the theory forces me to passively watch five other victims die.

\subsection{Non-Absolute Deontology (Ellis’ Paper)}
Under traditional Deontology, moral rules are absolute. For example, you should never kill under any circumstance.
This leads to Moral Catastrophes.

\noindent Imagine this scenario: a hidden nuclear device is about to explode. Unless you (Person A) violate your duty and torture an innocent person (Person B), ten, a thousand, or a million innocent people will die in the blast.
Since traditional Deontology forces us to let the millions die, we look to a softer theory called Non-absolute Deontology or Threshold Deontology.

\noindent\textbf{Non-absolute Deontology} says that moral rules are absolute until the consequences of following them become completely disastrous.
Once the threshold is crossed, the absolute rule breaks, and you are allowed to look at the consequences (and do the "bad" thing).

Anthony Ellis attacks this idea of a threshold. He argued that no plausible version of deontology can successfully capture what morality requires of us.
\textbf{Hijacked Plane Example}: Imagine terrorists hijack a plane with 400 people on board, and the terrorists threaten to blow up the plane unless you murder someone.
Ellis breaks the choice down into a mathematical spectrum:
\begin{itemize}
    \item Act 1: Murdering 1 person to save 1.
    \item Act 2: Murdering 1 person to save 2.
    \item ...all the way to Act 400: Murdering 1 person to save 400.
\end{itemize}
A non-absolutist deontologist would say that it is wrong to murder unless the consequences of refraining are sufficiently bad. So, they would easily say that Act 1 is wrong, but Act 400 is right.
Ellis' counter is: At exactly what point does the murder become right?
Any threshold point chosen is completely arbitrary. (If you say 399 people, then why not 398 people?)

\subsection{Degree of Wrongness}
Non-absolute deontologists defend against Ellis by introducing the concept of Degree of Wrongness.
They argue that the threshold isn't arbitrary because it scales based on Degrees of Wrongness. Some actions are "more wrong" than others, which means they require a much higher threshold of catastrophe to justify breaking them.
\begin{itemize}
    \item Stealing a car has a relatively low degree of wrongness. You might be morally allowed to steal a car to drive someone to the hospital and save one life.
    \item Murder has a massive degree of wrongness. You might need to save a million lives to cross the threshold and justify a murder.
\end{itemize}
They try to prove that the threshold system has a logical, sliding scale rather than just a random cutoff point.

\vspace{1em}
\noindent\textbf{Legal Model Trap (Ellis's Critique)}
To justify "Degrees of Wrongness," some deontologists try to borrow the structure of the legal system (which categorizes crimes into degrees, like 1st-degree vs 2nd-degree murder).

\noindent Ellis's Objection: This fails because the legal system is a man-made, conventional construct. Legal cutoffs are inherently arbitrary decisions made by judges/politicians.

\noindent Deontology is based on Intrinsic Wrongness (objective, universal truths based on reason). You cannot use a man-made, arbitrary legal framework to solve a problem about absolute, intrinsic moral facts.
Doing so reduces universal morality to a mere human convention.

\subsection{Paradox of Relative Stringency}
Ellis attacks the idea of Degrees of Wrongness, by arguing that they triggered the Paradox of Relative Stringency.
The entire foundation of Deontology is that moral rules are absolute duties (Categorical Imperatives).
By definition, an absolute cannot have "degrees." You cannot have one absolute rule that is somehow "stricter" or "more absolute" than another absolute rule.
Further, if you put weights on rules and comparing those weights against the weights of the consequences, Deontology turns into Utilitarianism.

\noindent\textbf{Conclusion}: Non-absolute Deontology is a failed theory.

\newpage

\section{Virtue Ethics}
Aristotle's Virtue Ethics is a pluralistic framework that contrasts with action-focused, monistic theories like Utilitarianism and Deontology. Instead of asking "What should I do?", Virtue Ethics asks "What sort of person should I be?".
\subsection{The landscape of normative ethics}
We have seen Utilitarianism and Kantian Deontology. They are classified as "monistic theories",
meaning they rely on a single, overarching principle to determine if an action is right or wrong. Both of these theories ask the fundamental question: "What should I do?".
\begin{itemize}
    \item Utilitarianism: This framework dictates that an action is morally right if it maximizes overall well-being, and wrong if it does not . This implies there is usually only one strictly "right" action in any given scenario.
    \item Deontology (Kantian): This framework argues that morality is not just about consequences. An action is morally right if and only if it can be universally applied as a rule, and if it treats human beings as ends-in-themselves rather than mere means to an end.
\end{itemize}
Virtue ethics is a "pluralistic" theory. This means it acknowledges that moral situations are complex, and it is entirely possible for there to be two different, equally right actions at the same time. Furthermore, Virtue Ethics shifts the foundational question of morality entirely.
It asks "What sort of person should I be?", "What is noble?", and "What is the chief good for humans?".
\subsection{The core of Virtue Ethics}
This shift in questioning represents a move from deontic claims to axiological claims.
\begin{itemize}
    \item Deontic claims evaluate actions (e.g., determining if an action is permissible, wrong, or right).
    \item Axiological claims evaluate things, events, states of affairs, and people—specifically their character traits, emotions, and intentions (e.g., determining if a person is virtuous, vicious, praiseworthy, or blameworthy) .
\end{itemize}
According to this framework, an action is morally right if and only if it is what a virtuous person would do in that specific situation.
\subsection{Defining Virtues and Moral Exemplars}
To figure out what a "virtuous person" would do, we must understand what a virtue actually is.\newline
A virtue is a character trait that manifests as a deep commitment to think, act, and feel in specific ways.\newline

\smallskip
\noindent Some virtues are specific to a role—for example, it is good for a doctor to possess the virtue of medical skill and knowledge. However, moral virtues (like courage, honesty, fairness, benevolence, and patience) are not specific to any profession.
They are character traits that are considered good for any human being to possess because they are necessary for living a "good human life".\newline
Because Virtue Ethics doesn't provide a strict rulebook, we rely on "moral exemplars" to guide our decisions. A moral exemplar is an ideally virtuous role model whose character makes them worthy of imitation.\newline
To make a moral decision using this theory, you simply insert a moral exemplar's name into this prompt: "In this situation, what would [Name] do?".
A crucial detail here is that a moral exemplar does not need to be a perfect human being who perfectly embodies every single virtue; they might even have vices. You are only meant to imitate them in situations where the specific virtue they excel at is relevant.\newline
For example, you should imitate Martin Luther King Jr. when you are faced with bigotry and injustice. However, because he plagiarized parts of his PhD dissertation, he would not be the correct moral exemplar to imitate when you are writing an academic paper.

\subsection{Aristotle's framework}
\subsubsection{Context: A Purpose-Driven View of Life}
To understand Aristotle's framework, we have to look at his view of the world. Aristotle believed that everything in the universe has a specific purpose or function. For example, the function of a knife is to cut, and a "good" knife is one that cuts well.
He applies this same logic to human beings through his Function Theory.

\smallskip
\noindent According to Aristotle, while we share basic functions like growth with plants and perception with animals, the function peculiar to humans is the "rational principle"—our ability to reason.
Therefore, a "good" human is one who reasons well, and virtues are the specific character dispositions that enable a person to function well and be good .

\subsubsection{Eudaimonia: The Ultimate Goal}
If we function well and live virtuously, what is the end result? Aristotle calls this Eudaimonia, which translates to "The Human Good," "Happiness," or "Living Well".
\begin{itemize}
    \item The nature of Eudaimonia: It is the ultimate goal of human existence because it is pursued entirely for its own sake and is considered more valuable than any single activity.
    \item The requirements: Achieving Eudaimonia is not a temporary state of mind; it requires using reason well over the course of an entire lifetime.
    \item The role of luck: Interestingly, Aristotle acknowledges that being virtuous isn't always enough on its own. Living well also requires "good fortune" and the possession of external goods like wealth, friends, and power.
\end{itemize}

\subsubsection{How Do We Become Virtuous?}
Aristotle divides virtue (excellence) into two distinct categories, which are acquired in different ways:
\begin{enumerate}
    \item Intellectual Virtues: These are acquired through teaching and learning.
    \item Moral Virtues: These are acquired through habit.
\end{enumerate}
Aristotle argues that moral virtues do not arise in us by nature. You are not born courageous or honest. Instead, virtue is a predisposition built through repeated practice. You become brave by repeatedly doing brave things until it becomes second nature.

\subsubsection{The Doctrine of the Mean}
When faced with a situation, how do we know what the virtuous action actually is? Aristotle provides a framework called the Doctrine of the Mean.
\begin{itemize}
    \item The Golden Mean: Virtue is defined as a state of balance between two extremes. Both of these extremes—the "Excess" and the "Deficiency"—are considered vices.
    \item Relative to us: This mean is not a strict mathematical average. It is "relative to us" and must be determined by reason and practical wisdom based on the specific situation. You must consider who you are, who is involved, and the specific facts of the scenario.
\end{itemize}
\textbf{Examples}:
\begin{itemize}
    \item Passions (Fear and Confidence): The excess is Rashness, the deficiency is Cowardice, and the mean virtue is Courage.
    \item Actions (Giving Money): The excess is Prodigality/Vulgarity, the deficiency is Meanness/Pettiness, and the mean virtue is Liberality/Magnificence.
\end{itemize}

\subsubsection{Virtue, Emotion, and Pleasure}
A vital part of Aristotle's framework is that virtue is not just about robotic action; it heavily involves our internal state. To be truly virtuous, one must be disposed to do the right thing, in the right way, at the right time, to the right person, for the right reason, and with the right amount of emotion.

\smallskip
\noindent The Ultimate Test of Virtue: To test if someone truly possesses a virtue, you must ask: Do they take pleasure in performing the action?. Through proper habituation, our relationship with pleasure is redefined so that doing the right thing actually feels good. For instance, a temperate person finds the act of abstaining from bodily pleasures to be pleasant in itself.

\subsubsection{The Contrast with Kant:}
This creates a sharp contrast with Kant's Deontology. Kant believed that moral worth is derived purely from duty, and that an action has no moral content if the person simply enjoys acting that way or does it out of natural sympathy.
For Kant, the highest moral worth is seen in someone who is deeply sorrowful and lacks all natural sympathy, but forces themselves to act charitably out of pure duty anyway.
For Aristotle, such a person is merely "continent" (exhibiting self-control) but not yet fully, joyfully virtuous.

\subsection{Advantages and Objections}
We have to evaluate it as a functional moral theory. Is it actually useful for guiding our lives?
\subsubsection{Advantages}
One of the most attractive features of Virtue Ethics is its realism; it aligns with our everyday experience that morality is deeply complex.
\begin{enumerate}
    \item No absolute rules: Unlike Kant's rigid rules, Virtue Ethics insists there is no single procedure that can dictate the precise right action in every single situation. While it uses virtue-based rules (like "Be generous"), it recognizes these are not absolute and do not apply uniformly to all cases.
    \item Striking a balance: The theory acknowledges that different virtues often compete. For instance, virtues requiring narrow concern (like loyalty to a family member) often conflict with virtues requiring broad concern (like general benevolence to society).
    \item The role of Practical Wisdom: To resolve these conflicts, a virtuous person uses "practical wisdom" and rational reflection to judge on a case-by-case basis which virtue should take priority.\newline
    Implication/Example: The slides contrast two scenarios to prove this point. In the story of Andrew and Belle, loyalty rightfully trumps benevolence. However, in the true story of Ted and David Kaczynski (the Unabomber), David's benevolence to society rightfully trumped his loyalty to his brother, leading him to turn Ted in.
    \begin{itemize}
        \item Scenario 1: Loyalty Trumps Benevolence (Andrew and Belle). Think of a common, lower-stakes ethical dilemma where a family member makes a mistake or needs your help in a way that slightly inconveniences the broader community. A purely utilitarian approach might demand you prioritize the "greater good" of society at all times. However, Virtue Ethics recognizes that being a good human means honoring the deep, narrow bonds of loyalty to the people closest to you. In the case of Andrew and Belle, the ethical choice is to prioritize loyalty over general benevolence.
        \item Scenario 2: Benevolence Trumps Loyalty (The Kaczynski Brothers). Ted Kaczynski was the infamous "Unabomber," who conducted a 17-year domestic terror campaign using mail bombs. His brother, David, eventually recognized Ted's distinct writing style in a published manifesto and realized his own brother was the killer.\newline
        David faced a devastating moral dilemma with two heavily competing virtues:
        \begin{itemize}
            \item Loyalty (Narrow Concern): Protect his brother from the FBI and the potential of the death penalty.
            \item Benevolence (Broad Concern): Protect the general public from future bombings and save lives.
        \end{itemize}
        In this extreme situation, the narrow virtue of familial loyalty was outweighed by the massive harm being done to society. David chose to turn his brother in. The slides use this to demonstrate that in contexts involving severe harm, broad benevolence rightfully trumps narrow loyalty.
    \end{itemize}
\end{enumerate}
\subsubsection{Objections}
\begin{enumerate}
    \item Conflicting Advice: If we are supposed to ask "What would a virtuous person do?", what happens when different virtuous people give conflicting advice? This suggests the theory cannot give us a clear answer.
    \begin{itemize}
        \item The Response: Virtue ethicists argue that the advice of perfectly virtuous people would not actually conflict. Because perfectly virtuous individuals possess practical wisdom, they know exactly how to rationally balance conflicting considerations.\newline
        The implication: If you are looking at less-than-perfect role models whose advice does conflict, the theory advises you to follow the person who best exemplifies the specific virtue that is most relevant to the situation at hand.
    \end{itemize}
    \item What is the good life?: Virtue Ethics claims moral virtues are absolutely necessary for a good human life. Critics argue this is false; it seems entirely possible to live a "good," successful, or happy life without being a morally virtuous person.
    \begin{itemize}
        \item The Response: Aristotle's defense rests on his Function Theory: a good life inherently requires the successful exercise of reason. To reason successfully, one needs practical wisdom, and someone with practical wisdom will naturally acquire the moral virtues.\newline
        Context: Modern psychology also supports this Aristotelian view, showing a strong connection between virtue and well-being. For example, research indicates that people who act benevolently are generally happier.
    \end{itemize}
    \item Are character traits even real? (Situationism): This is a modern critique driven by psychological research, specifically "Situationism." Famous studies like the Stanford prison experiment (1971) and the Milgram electric shock experiment (1963) show ordinary people committing shocking atrocities when placed in specific environments.
    This suggests that external factors (the situation), rather than internal character traits (virtues), dictate how we act.\newline
    Responses:
    \begin{enumerate}
        \item Lack of virtue: The people who committed atrocities in these experiments simply didn't possess the relevant virtues in the first place.
        \item A continuum: Being virtuous isn't a strict "all-or-nothing" binary; it exists on a continuum, so being virtuous in some situations still has moral value.
        \item Frail traits: Philosopher Robert Adams (2006) suggests that while human character traits might be "frail and fragmentary," they still count as virtues.
    \end{enumerate}
\end{enumerate}

\subsubsection{The Perennial Problem}
\textbf{The Objection}: Philosopher Julia Annas (2004) outlines a fundamental problem: If right action is defined as what a virtuous person would reliably do, how do we actually identify a truly virtuous person? Furthermore, if the "virtuous person" is just a theoretical ideal and doesn't exist in reality, how does the theory actually work to give us actionable advice?\newline
\textbf{The Response (Annas's Solution)}: Annas argues that critics are misunderstanding the purpose of the theory. Virtue Ethics is not a "technical manual" designed to merely tell you what to do in every scenario. Instead, it is a developmental framework. It is designed to help you understand your own moral reasoning and justify your choices, because ultimately, your decisions are your own.

\subsection{Application - Corporate Whistleblowing}
A whistleblower is an employee (an "insider") who informs a government agency or the general public about an illegal, harmful, or unethical activity being committed by their own business or institution.\newline
The slides point to Edward Snowden, who leaked highly classified information from the National Security Agency, as a famous example of a whistleblower.
\subsubsection{The dilemma - Competing Virtues}
Whistleblowing creates a profound moral dilemma because there are highly compelling reasons both for and against taking action.
\begin{itemize}
    \item Reasons For: It exposes the truth and can prevent harmful or unethical activity by the business.
    \item Reasons Against: It requires acting directly against your employer, often comes at a massive personal and professional cost, and might not even be successful in changing the bad behavior.
\end{itemize}
When viewed through the lens of Virtue Ethics, this dilemma is a clash of competing character traits. A potential whistleblower is forced to balance the virtues of Justice and Compassion (which demand exposing the harm to protect the public) against the virtues of Loyalty, Obedience, and Prudence (which demand protecting the company, honoring contracts, and protecting oneself).
\subsubsection{The Virtuous Response and Practical Wisdom}
If you were to ask, "What would a moral exemplar do in this situation?", the answer is notoriously difficult to pin down.
\begin{itemize}
    \item No Simple Rule: Because Virtue Ethics rejects absolute rules, there is no single procedure to follow. The moral exemplar must rely entirely on "practical wisdom" to weigh exactly how much personal sacrifice is appropriate to prevent a specific harm.
    \item Multiple Right Answers: While some virtue ethicists argue there is still only one truly right answer for any specific scenario, others argue that in highly complex situations, both whistleblowing and keeping silent might be morally permissible. The dilemma can become so convoluted that even a perfectly virtuous person might be genuinely undecided on the best course of action.
\end{itemize}

\subsubsection{DeGeorge's Model of Whistleblowing}
To provide a more structured way to navigate this, the slides introduce a model created by moral philosopher Richard DeGeorge. Even though DeGeorge wasn't strictly a virtue ethicist, his model illustrates how we can rationally balance competing virtues to determine when whistleblowing is justified.\newline
\textbf{When is whistleblowing morally permitted?} It is permitted when an employee successfully balances Justice and Compassion with Loyalty and Obedience. According to the model, this occurs when:
\begin{itemize}
    \item The firm's product or policy will do serious harm to employees, users, or the public.
    \item The employee has already reported the concern to their immediate supervisor.
    \item All internal possibilities to fix the issue have been completely exhausted, yet the concern remains unaddressed.
\end{itemize}
\textbf{When is whistleblowing morally required?} The action elevates from merely permitted to strictly required when the employee balances Justice and Compassion with Prudence (acting with careful, rational foresight). This occurs when the previous conditions are met, plus:
\begin{itemize}
    \item The employee possesses convincing, documented evidence that serious harm will occur.
    \item The employee has good reason to believe that publicizing the information will actually bring about the required changes, ensuring their sacrifice isn't in vain.
\end{itemize}


\newpage

\section{Questions}
\subsection{Metaethics}
\begin{enumerate}
    \item Both non-cognitivists and error theorists agree that morality is an illusion because it does not exist. TRUE or FALSE.
    \begin{answer}\textbf{FALSE.} Non-cognitivism only denies that moral statements are truth-apt (they express emotion), and being a non-cognitivist does not make one a moral skeptic. Error theory says moral statements \emph{attempt} to state facts but are all false. It is Extreme Nihilism (an extension of error theory), not non-cognitivism, that calls morality an illusion --- so the two views do not share that claim.\end{answer}
    \item Cognitivists think that when moral judgments are true, they are always proven true in virtue of God’s will. TRUE or FALSE?
    \begin{answer}\textbf{FALSE.} Cognitivism only claims that moral statements are truth-apt; \emph{what} makes them true varies (e.g.\ natural facts under ethical naturalism). Grounding moral truth in God’s will is specifically Divine Command Theory / ethical supernaturalism, not a commitment of all cognitivists.\end{answer}
    \item According to Harman, we cannot use evidence in ethical observations to justify the truth of moral principles because moral principles are never true. TRUE or FALSE?
    \begin{answer}\textbf{FALSE.} Harman’s actual reason is that moral facts are explanatorily idle --- unnecessary to explain our observations (Ockham’s Razor), so we have no evidence for them. He does not claim moral principles are ``never true'' (that is error theory). The stated reasoning is wrong.\end{answer}
\end{enumerate}
\subsection{Moral Relativism}
\begin{enumerate}
    \item If you believe in (empirical) cultural relativism, then you must be a cognitivist about moral judgments. TRUE or FALSE?
    \begin{answer}\textbf{FALSE.} Descriptive (empirical) cultural relativism is just the observation that beliefs about right and wrong vary across cultures. It says nothing about whether moral judgments are truth-apt, so it carries no cognitivist commitment.\end{answer}
    \item Utilitarians are not moral relativists because…
    \begin{enumerate}
        \item The utility principle applies to everyone regardless of their culture.
        \item Benthamic utilitarians are often interpreted as ethical naturalists.
        \item There can be only one universally right action for every moral predicament.
        \item All of the above.
        \item a and c only.
    \end{enumerate}
    \begin{answer}\textbf{(d) All of the above.} The utility principle is universal and impartial (a); Benthamic utilitarians are ethical naturalists, who hold that objective moral facts exist --- the opposite of relativism (b); and it posits a single universal standard that fixes the right action (c). All three make utilitarianism universalist rather than relativist.\end{answer}
\end{enumerate}

\subsection{Utilitarianism}
\begin{enumerate}
    \item The principle of utility requires that we always perform the action that maximizes well-being, therefore, it entails that we ought to always choose the action that benefits the greatest number of people. TRUE or FALSE?
    \begin{answer}\textbf{FALSE.} The utility principle maximizes \emph{total net utility}, not head-count. A huge benefit to two people can outweigh a tiny benefit to a hundred, so ``the greatest number of people'' does not follow.\end{answer}
    \item The following is an example of an agent-relative reason
    \begin{enumerate}
        \item Every parent ought to save their own children because parents have special obligation to their children.
        \item Everyone has a duty to save a drowning child when it is a very easy rescue.
        \item Everyone has a duty to save people when they themselves are feeling happy.
        \item Every medical doctor ought not to kill an innocent patient even if he can save more people because medical doctors have made promises to do no harm.
        \item Options a, b and d
        \item Options a and d
    \end{enumerate}
    \begin{answer}\textbf{(f) Options a and d.} Both are universal rules applied to a specific relationship or role (a parent to their own child; a doctor bound by their promise), which is what makes a reason agent-relative. (b) is an agent-neutral duty that holds for everyone; (c) is a subjective feeling, not an agent-relative moral rule.\end{answer}
    \item Utilitarianism is a defective theory because we cannot always predict the consequences of our actions, so it cannot help us deliberate and make decisions about what to do. (Note: True/False based on how utilitarians answer the Deliberation objection)
    \begin{answer}\textbf{FALSE.} Utilitarians answer the Deliberation objection by separating the \emph{Standard of Rightness} (what makes an act correct) from the \emph{Decision Procedure} (how we actually choose --- proven rules of thumb, the basketball analogy). So the objection does not show the theory is defective.\end{answer}
\end{enumerate}

\subsection{Deontology}
\begin{enumerate}
    \item Kantian deontologists argue that we have a perfect duty to help others. TRUE or FALSE?
    \begin{answer}\textbf{FALSE.} Helping others is an \emph{imperfect} duty (flexible in when, how, and toward whom it is fulfilled), not a perfect one.\end{answer}
    \item Under Kant's Principle of Humanity, it is always immoral to use another human being as a means to an end. TRUE or FALSE?
    \begin{answer}\textbf{FALSE.} It is immoral only to use someone as a \emph{mere} means. Using another as a means is permissible when they share the purpose (the barista example). The word ``mere'' is the whole point.\end{answer}
    \item According to Kant, imperfect duties are less important than perfect duties because imperfect duties are merely optional suggestions, not actual moral obligations. TRUE or FALSE?
    \begin{answer}\textbf{FALSE.} Imperfect duties are still strict moral requirements: you \emph{must} adopt the goal (e.g.\ helping others). Only the when, how, and toward whom is left to your discretion --- they are not optional suggestions.\end{answer}
    \item According to the principle of humanity, we ought to treat humans never as mere means, but always as ends in themselves. This means that it is morally wrong to be the CEO of any company because I will be making people work for me. TRUE or FALSE?
    \begin{answer}\textbf{FALSE.} Employees share the purpose of the arrangement (like the barista getting paid for the transaction), so they are not being used as a \emph{mere} means. Employing people is not a violation of the Principle of Humanity.\end{answer}
    \item Based on Ellis’ paper, deontologists can appeal to legal prohibitions to explain the idea of “degree wrongness” but this can be problematic because…
    \begin{enumerate}
        \item The idea of intrinsic wrongness is essential to the deontologists.
        \item Even if we accept the very idea of degree of wrongness, we still cannot justify the cut-off point without being arbitrary.
        \item It is not the case killing is always the worst wrongdoing even though we tend to take the duty to not kill much more seriously than other duties.
        \item All of the above.
        \item b and c only.
    \end{enumerate}
    \begin{answer}\textbf{(d) All of the above.} Ellis’ Legal Model Trap: deontology rests on \emph{intrinsic}, reason-based wrongness (a), but legal categories are a man-made convention whose cut-offs are arbitrary (b) and whose severity does not track intrinsic moral wrongness (c). Borrowing the legal framework therefore reduces universal morality to human convention.\end{answer}
\end{enumerate}

\end{document}